\chapter{熱流束回帰と沸騰開始判定の手法}\label{ch:methods}
\section{音響の周波数解析とスペクトログラムの作成}
\subsection{高域通過と時間周波数表現}
本研究では，保存版\texttt{waterflow\_20260817\_1s}の線形power配列を用いる．生成仕様\cite{localData}は，通過域端500\,Hz，阻止域端400\,Hz，通過域損失1\,dB，阻止域減衰40\,dBの高域通過処理を含む．この前処理により抑制された低周波の有用性を，現在の3\,kHz入力の結果だけから判断することはできない．また，通過域端500\,Hzは完全な矩形除去の境界ではなく，フィルタの仕様値である．

1秒の離散音響信号$x[n]$に対し，窓$w[n]$を移動しながら短時間フーリエ変換を行う．窓長を$N_w$，フレーム間隔を$H$とすると，
\begin{equation}
 X(t,k)=\sum_{n=0}^{N_w-1} x[n+tH]w[n]\exp\left(-j\frac{2\pi kn}{N_w}\right)
\end{equation}
と表せる．実装はHann窓とその振幅補正を用い，FFTの片側振幅を二乗した値を保存する．従って，本稿でSTFT powerと呼ぶ配列と，別の診断で用いたWelch法によるpower spectral density（PSD）は，数値尺度を区別する．校正音圧に基づく絶対PSDと同一視しない．

生成処理では672フレーム，FFT窓長1,344点を用い，1秒44,100点に対してフレーム間隔は63点となる．元のFFT周波数間隔は約32.81\,Hzである．上限3\,kHzに対応する周波数範囲を切り出し，時間・周波数の2軸を$224\times224$へresizeし，float32で保存する．resizeにより画素数が増えても，元の周波数分解能が高くなるわけではない．後述の34要約特徴で用いるHz対応は，このresize後配列に対する近似座標である．

\subsection{モデルへの共通入力}
各モデルへ渡す基礎配列を$P_i\in\mathbb{R}_{\ge0}^{224\times224\times1}$とする．第1軸は1秒内の時間フレーム，第2軸は周波数，第3軸は1チャンネルである．CNN系では，
\begin{equation}\label{eq:logpower}
 L_i=\log\left(1+\frac{\max(P_i,0)}{10^{-12}}\right)
\end{equation}
をモデル内部で計算する．これは区間ごとの最大値を1にする正規化とは異なり，powerの絶対的な大小関係を残した対数圧縮である．RF系は配列の圧縮表現を，追加2モデルは同じ配列から計算した固定要約特徴を用いる．従って，同じ基礎音響を用いる一方で，モデルごとに入力表現が異なる．

\section{回帰モデルへの入力特徴}
\subsection{PCAによる圧縮表現}
元RFでは，power配列を平坦化し，学習側で求めた主成分分析（PCA）で100次元へ圧縮する．PCAの中心と主成分は各fit側だけから求め，held-out側と評価側へ固定適用する．保存・再読込時の数値的な差による木の分岐変化を抑えるため，実装では主成分行列の配列配置を揃え，変換後を小数点以下6桁へ丸める．この処理は再現のための数値条件であり，音響的な重要度を定義するものではない．

\subsection{34個の音響要約特徴}
追加2モデルは，学習パラメータを持たない同じ34特徴変換を使用する．各特徴の内訳は表\ref{tab:features}のとおりである．
\begin{table}[htbp]
\centering\small
\caption{34要約特徴の構成．帯域値はresize後の近似周波数座標に基づく．}
\label{tab:features}
\begin{tabularx}{\textwidth}{L{36mm}rY}
\toprule
特徴群 & 個数 & 定義・対象 \\
\midrule
帯域log-power & 4 & 512--1,000，1,000--2,000，2,000--3,000，2,100--2,500\,Hzの配列平均の常用対数 \\
全域log-power & 1 & 配列全体の平均の常用対数 \\
帯域log比 & 2 & 1,000--2,000\,Hzに対する2,000--3,000および2,100--2,500\,Hzの比の常用対数 \\
周波数形状 & 2 & 周波数重心，正規化スペクトルentropy \\
時間変動 & 1 & 時間フレームごとの平均powerの変動係数 \\
相対log-power & 24 & 0--3,000\,Hzを125\,Hz刻みにした各帯域のlog-powerから全域log-powerを引く \\
\bottomrule
\end{tabularx}
\source{特徴変換実装\cite{localCode}．}
\end{table}

帯域$B$に対応する画素集合を$\mathcal{K}_B$とし，時間平均も含む帯域平均を
\begin{equation}
 p_B=\frac{1}{224|\mathcal{K}_B|}\sum_{t=1}^{224}\sum_{k\in\mathcal{K}_B}P(t,k)
\end{equation}
とする．log-powerには$\log_{10}\max(p_B,10^{-20})$を用いる．帯域画素の端点は$\operatorname{round}(224 f/3000)$で与える近似であり，理想的な帯域積分としては扱わない．重心とentropyでは，時間平均スペクトルを総和で正規化した重みを用いる．entropyは$\log224$で規格化する．

時間power変動係数は，$a_t=224^{-1}\sum_kP(t,k)$の標準偏差を平均で除したものである．これは変動の大きさを表すが，時間の順番を保持しない．他の特徴も時間集約を用いるため，1秒内の時間行を並べ替えても34特徴は不変となる．「時間に関する情報を全く含まない」表現ではなく，「時間順序を捨てる」表現と位置づける．日付，熱流束ラベル，正解の雑音レベル，無雑音参照は特徴へ含めない．

\section{熱流束回帰モデル}
\subsection{従来3モデル}
本稿で元3と呼ぶのは，RF，CNN＋Transformer，AlexNetである．RFは一般的な呼称を引き継ぐが，実体はPCAとXGBoostの\texttt{XGBRFRegressor}\cite{xgbrf}であり，scikit-learnのRandomForestRegressorと同一の実装ではない．主比較では木数100，最大深さ12，subsampleおよびcolsample\_bynodeを0.6としている．

CNN＋Transformerは，式\eqref{eq:logpower}の後に5つの畳み込み層とpoolingを配置し，抽出特徴を時間軸のtoken列へ変換する．各tokenを64次元へ射影し，位置埋め込みを与え，4 head，key次元16，中間次元256のTransformer encoderを2段重ねる．最後に時間平均と線形回帰出力を用いる．attentionの考え方はTransformer\cite{vaswani2017}に基づくが，過去の保存キー\texttt{conformer}だけから標準Conformerの再現と説明してはならない．本稿では実装に即した名称を用いる．

AlexNet系\cite{krizhevsky2012}では，対数圧縮後の入力へ5つの畳み込み層を適用し，Flattenと4,096ユニットの全結合層2段から線形出力を得る．CNN＋Transformerと前段の構造が近いため，名称の違いだけで完全に独立な特徴を利用しているとは考えない．両モデルの依存差は，第5章の入力対照で確かめる．

\subsection{追加2モデル}
HistGradientBoosting（HGB）は，34特徴に対する二乗誤差を逐次的に補正する木モデルである．histogramによる分割を用いる実装\cite{sklearnHGB}とし，最大反復150，学習率0.05，最大葉数31，最小葉標本数20，L2正則化1.0，early stoppingなしとする．要約特徴から非線形な関係を学習し，元3が残す誤差へ異なる予測を与えるかを調べる．

ExtraTrees（ET）は，特徴・分割位置のランダム化を用いる木の平均\cite{geurts2006}である．同じ34特徴に対し256本の木，最小葉標本数2，max\_features 1.0，bootstrapなしを用いる．HGBと入力を揃えることで，逐次補正とランダム化した木の平均という学習方式の違いを比較できる．ただし，元3との比較では表現と学習器を同時に変えるため，性能差を入力表現だけの効果と断定しない．

追加2モデルは既存の学習側比較を経て採用された条件\cite{localInput}であり，本稿の執筆にあたって候補探索を追加していない．採用理由は，少数要約特徴を用いて異なる予測を得るという仮説と，保存された比較結果の双方から説明する．モデル数が増えたこと自体は効果の証拠ではない．

\section{学習・評価方法}
第2章の1,620学習区間内をshuffle chunk 3-foldに分け，二つのfoldで学習して残りのfoldを予測する処理を3回行う．これにより，全学習区間に対するout-of-fold（OOF）予測を一つずつ得る．OOFは当該chunkをfitに用いない予測であるが，同じWAVの別chunkがfit側に入る．従って，OOFも未知WAV評価ではない．PCAおよび目的変数のscalerは，各fit側のみに基づく．

統合重みをOOFから決めた後，1,620区間全体で最終モデルを学習し，540評価区間へ固定適用する．CNN系は内部3-foldと最終fitを含む計8学習で200 epochs完走した記録がある．学習率はCNN＋Transformerが0.001，AlexNetが0.003，batch sizeはそれぞれ12，8である\cite{localMainRun}．目的変数のMinMax変換は学習側で求め，予測を逆変換して保存する．保存予測の単位は\wm{}，主表では1,000で除して\kw{}とする．

全モデルのハイパーパラメータが十分に最適化されたとは主張しない．また，training lossの減少や200 epochsの完走は，未知区間上の最良epochや過学習がないことの証拠ではない．本稿は保存された固定条件の比較を扱う．

評価区間は当該fitへ入れていないが，研究の途中で結果を繰り返し参照し，モデル構成や方針を検討している．CawleyとTalbot\cite{cawley2010}が論じるモデル選択と評価の関係を踏まえ，本稿の540区間を完全に未参照の最終試験とは呼ばない．本稿の数値は研究開発過程での対応比較として解釈する．
\need{V01}{200 epochsが最良であるという主張を採る場合に限り，学習側で独立性を定めたvalidation曲線とepoch選択を要する．本稿の固定条件の比較を記述するために，追加学習を自動的な必須作業とはしない．}

\section{アンサンブル学習による予測の統合}
\subsection{元3の重みと固定ブロック}
モデル$m$のOOF予測を$z_{im}$，正解を$q_i$とする．元3 MSE統合は，
\begin{equation}\label{eq:core}
 \min_{\bm a}\frac{1}{N_{\rm tr}}\sum_i\left(\sum_{m=1}^{3}a_mz_{im}-q_i\right)^2,
 \quad\sum_ma_m=1,\quad a_m\ge10^{-6}
\end{equation}
で内部比率を求める．従来performance統合は，各モデルのOOF MSEを正解分散で割った誤差に対する逆数を規格化する方式である．元3 MSEは予測を組み合わせた誤差を直接最小化し，performanceは個別誤差から重みを与えるという違いがある．いずれも評価側の正解を重みfitに使用しない．

元3の重み付き予測を$z_{i0}$という一つのブロックとして固定する．5モデルではHGBとETを加え，
\begin{equation}\label{eq:five}
 \hat q_i=b_0z_{i0}+b_Hz_{iH}+b_Ez_{iE}
\end{equation}
とする．係数はclean OOFの全域MSEを最小にするように求め，
\begin{equation}
 b_0+b_H+b_E=1,\qquad b_0\ge0.25,\quad b_H\ge0.05,\quad b_E\ge0.05
\end{equation}
を課す．4モデルでは対応する追加モデル一つだけを含める．この下限は元3と追加モデルを比較構成内に保持するための設計条件であり，25\%や5\%が物理的・普遍的な最適値であることを意味しない．最適化はSLSQPで行い，制約と収束を確認して重みを保存する．

\subsection{比較構成と保存重み}
5単体に加え，元3 performance，元3 MSE，元3＋HGB，元3＋ET，5モデル統合，5等重みの計6統合を比較する．単体と統合は同じ基礎予測を共有し，違いは保存された重みである．主比較の重みを表\ref{tab:weights}に示す．
\begin{table}[htbp]
\centering\small
\caption{全7条件で固定する個別モデルの重み [\%]．表示丸めにより総和が100とわずかに異なる場合がある．}
\label{tab:weights}
% Generated from saved evidence; do not hand-edit.
\begin{tabular}{lrrrrr}
\toprule
方式 & RF & CNN+Tf & AlexNet & HGB & ET \\
\midrule
3-perf & 16.94 & 55.15 & 27.91 & 0.00 & 0.00 \\
3-MSE & 5.85 & 71.32 & 22.83 & 0.00 & 0.00 \\
3+HGB & 2.08 & 25.38 & 8.12 & 64.42 & 0.00 \\
3+ET & 1.46 & 17.83 & 5.71 & 0.00 & 75.00 \\
5-model & 1.46 & 17.83 & 5.71 & 27.25 & 47.75 \\
5-equal & 20.00 & 20.00 & 20.00 & 20.00 & 20.00 \\
\bottomrule
\end{tabular}

\end{table}

5モデルの元3ブロックは25\%，HGBは約27.25\%，ETは約47.75\%である．元3ブロックの下限に達していることは，無制約なら同じ構成が選ばれるという保証ではない．元3比率を自由にする対照は第5章で扱い，構成保持の意図と数値上の最適化を分けて考察する．

\section{雑音の付加方法と学習方針}
\subsection{固定基準による水流音付加}
水流音の元資料は\texttt{water\_flow\_125.wav}である．基準SNRを$s$\,dB，切り出した雑音区間を$n_i$，そのRMSを$R(n_i)$，共通基準RMSを$R_{\rm ref}$とすると，付加信号を
\begin{equation}\label{eq:noise}
 x_i^{(s)}=x_i+\frac{R_{\rm ref}10^{-s/20}}{R(n_i)}n_i
\end{equation}
とする．雑音区間を規格化するため，付加雑音のRMSは，元信号の熱流束や個々のRMSに依存しない．同じchunkには雑音レベル間で同じ水流音区間を対応させる．

保存された$R_{\rm ref}$は$0.0004822199855309985$であり，生成時に設定された3日分・49 WAVのフィルタ後RMSの中央値である\cite{localData}．主比較はそのうち6月の36 WAVを用いる．7月9日の録音は雑音基準の算出母集団に含まれるが，本稿の性能評価へ追加したものではない．この基準はデータ生成時の共通負荷を定義するものであり，実運用で未知録音から何も参照せず設定できた校正値としては主張しない．また，校正音圧レベルではなくdigital amplitude上の値である．

各chunkの実現SNRは，
\begin{equation}\label{eq:realized}
 s_i=20\log_{10}\frac{R(x_i)}{R_{\rm ref}10^{-s/20}}
 =s+20\log_{10}\frac{R(x_i)}{R_{\rm ref}}
\end{equation}
となり，$s$と一致するとは限らない．基準SNR $-20$\,dBは，基準に対して雑音RMSが10倍，powerが100倍となる設定であり，全区間の実現SNRが$-20$\,dBという意味ではない．本稿の図表では基準SNRと明記する．

旧生成の診断では，元音の強さに応じて付加雑音も変わり，付加処理が熱流束の手掛かりを持ち得る経路が確認された．現版の共通基準はこの経路を分けるための条件である．ただし，元信号の強さに依存する相対SNRそのものを，全ての用途で不正な設計と一般化しているわけではない．本研究の目的に対して，付加負荷を明示的に統制する方式を採る．

\subsection{固定モデル評価と条件別再学習の違い}
本稿の主系列はclean\_onlyであり，無雑音で学習した同じモデル，PCA，scaler，統合重みを全条件へ固定適用する．これは，学習時に含まれない水流音の付加に対して予測がどう変わるかという問いに答える．

一方，matchedは各雑音条件で再学習する方式であり，雑音条件を学習した場合の性能を問う．matched自体を不具合とはしない．旧入力生成の修正と，主評価として固定モデルを選ぶ判断は別である．また，単調に悪化する曲線が得られることを主系列の採用理由にはしない．

\section{予測熱流束を用いた沸騰開始判定}
収録日$d$の区間$i$に対し，日別ONBを$\qonb$として，
\begin{equation}\label{eq:decision}
 y_i=\mathbf{1}(q_i\ge\qonb),\qquad
 \hat y_i=\mathbf{1}(\hat q_i\ge\qonb)
\end{equation}
と定義する．閾値は実験側の採用値であり，評価側の見逃しを最小にするように後から選び直さない．式\eqref{eq:decision}は既知の日別閾値を用いる解析上の判定であり，未知の実験条件のONB閾値を音響だけから同定する問題とは異なる．

本研究は段階ごとに収録された区間を評価するため，真のONB発生時刻から何秒後に警報できたかを測定していない．「早期」の適合性は，開始段階の見逃しと，より高い熱流束まで進まないと全区間を陽性にできない程度から限定的に検討する．秒単位の遅延や実機での警報余裕とは区別する．

\section{予測性能の評価指標}
\subsection{回帰と領域別指標}
評価区間数を$N$とし，残差を$e_i=\hat q_i-q_i$とする．
\begin{align}
 \mathrm{RMSE}&=\sqrt{\frac{1}{N}\sum_i e_i^2},&
 \mathrm{MAE}&=\frac{1}{N}\sum_i|e_i|,\\
 \mathrm{bias}&=\frac{1}{N}\sum_i e_i,&
 R^2&=1-\frac{\sum_i e_i^2}{\sum_i(q_i-\bar q)^2}
\end{align}
とする．$R^2$は負値を取り得る．全域RMSEは540区間をまとめて求める集約値であり，fold別RMSEの単純平均とは区別する．近傍集合を$\mathcal N=\{i:|q_i-\qonb|\le0.1\qonb\}$とし，同じ式を集合内に適用する．ONB以上全体と，近傍を除いた高熱流束側を区別する．

雑音条件$s$の悪化量には，$D_m(s)=\mathrm{RMSE}_m(s)-\mathrm{RMSE}_m(\mathrm{clean})$を用いる．異なるモデル間の雑音時の誤差差と，同じモデルのcleanからの増加量を別々に示す．

\subsection{見逃し・誤報と順位評価}
式\eqref{eq:decision}の正解と予測からTP，FN，FP，TNを数える．ONB以上を陰性とする誤りがFN，ONB未満を陽性とする誤りがFPである．
\begin{equation}
 \mathrm{Recall}=\frac{TP}{TP+FN},\quad
 \mathrm{FPR}=\frac{FP}{FP+TN},\quad
 F_1=\frac{2TP}{2TP+FP+FN}
\end{equation}
とする．全ONB以上のRecallに加えて，記録上ONB段階だけの検知件数を示す．陽性数が多い高熱流束側で良い成績を得ることによって，開始段階の未達が見えにくくなるのを避けるためである．

ROC-AUCは予測熱流束の連続スコアに対する順位性能として補助的に扱う．日別閾値が異なる集合を生の予測値でまとめたAUCと，日別固定閾値の判定件数は同じ指標ではない．AUCが高くても採用閾値の近傍でFNが残るため，AUCから開始段階の検知達成を結論しない．

\subsection{検知到達熱流束}
日$d$に含まれる記録段階の集合を$\mathcal Q_d$とする．
\begin{equation}\label{eq:q100}
 q_{100,d}=\min\{q\in\mathcal Q_d:\ \forall i\ (d_i=d,\ q_i\ge q)\Rightarrow\hat y_i=1\}
\end{equation}
を定義する．該当する段階がなければ未到達とする．さらに$g_{100,d}=q_{100,d}-q_{\mathrm{ONB},d}$を示す．$q_{100}$は有限の評価集合について，それ以上の全区間が陽性となる最小の記録段階である．誤報ゼロは条件に含めない．常時陽性のモデルでも低い値を取り得るため，FN/FPと併記する．

この値を将来の100\%検知保証や，連続的な真のONB推定値と解釈しない．段階間隔にも依存し，1区間の予測変化で一段前進することがある．最大記録熱流束をCHFとした規格化も用いない．

\section{特徴利用と予測誤差の分析方法}
\subsection{領域分解と同方向・逆符号の比較}
元3と5の誤差差を，全域，ONB前，近傍，高熱流束側に分ける．保存予測での元3残差を$e_i$，5への調整を$\delta_i=\hat q_{5,i}-\hat q_{3,i}$とすると，
\begin{equation}\label{eq:decomp}
 \mathrm{MSE}_5-\mathrm{MSE}_3
 =\frac{2}{N}\sum_i e_i\delta_i+\frac{1}{N}\sum_i\delta_i^2
\end{equation}
が恒等的に成り立つ．調整が元の誤差と反対向きでも，調整量が大きすぎると悪化するため，符号と大きさの両方を見る．元3・HGB・ETの3予測群について同方向群と逆符号を含む群に分け，群内の二乗誤差減少の総和を全$N$で除して全体MSE減少への寄与を求める．これは予測値の損失分解であり，特徴や物理音源の因果寄与率ではない．

\subsection{共通入力加工と特徴除去}
既存の10月7日の対照\cite{localInput}では，同じraw power配列に時間行並べ替え，全域power 0.5倍・2倍，対象帯域power 0.1倍などの加工を与え，モデルを固定した予測変化を保存している．1秒内の時間配置，強度，周波数帯の依存を見る手順であり，熱流束の正解は変更しない．人工加工は元の音響分布から外れる場合があり，実際の水流音付加と同じ外乱ではない．

別の対照では，34特徴から絶対強度5特徴だけ，または相対形状等29特徴だけを残して学習し直している．固定モデルの入力加工と，情報を減らした後の再学習は異なる比較である．さらに，特徴変更と同時に統合重みを再fitすると二つの要因が混在するため，元の重みを固定する対照で分ける．第5章ではこれらの証拠を，誤差の説明，入力依存，物理的解釈の順に接続する．
